\chapter{Mathematics}\label{ch-math}

The arithmetic operators are in the language itself
(\hyperref[ch-numbers]{Numbers and Operators}). A few functions sit
beside them with no import, and the \verb|math| module holds the rest:
square roots, rounding, angles, logarithms and powers of $e$.

\section{The functions that need no import}

\begin{center}
\begin{tabular}{lP{8.5cm}}
\toprule
Written & Gives \\
\midrule
\verb|abs(x)| & \verb|x| without its sign \\
\verb|min(a, b, ...)|, \verb|min(xs)| & the smallest of the numbers, or of a list \\
\verb|max(a, b, ...)|, \verb|max(xs)| & the largest \\
\verb|sum(xs)| & the numbers of a list added up \\
\bottomrule
\end{tabular}
\end{center}

\begin{zeno}
print(abs(-3), abs(2.5), min(3, 1, 2), max([4, 9, 2]), sum([1.5, 2.5]))
\end{zeno}

\section{The math module}

\begin{center}
\begin{tabular}{lll}
\toprule
Written & Takes & Gives \\
\midrule
\verb|math.pi|, \verb|math.e| & & \verb|float| \\
\verb|math.sqrt(x)| & a number, 0 or more & \verb|float| \\
\verb|math.floor(x)|, \verb|ceil|, \verb|round| & a number & \verb|int| \\
\verb|math.sin(a)|, \verb|cos|, \verb|tan| & an angle & \verb|float| \\
\verb|math.asin(x)|, \verb|acos| & a number from $-1$ to 1 & \verb|rad| \\
\verb|math.atan(x)| & a number & \verb|rad| \\
\verb|math.atan2(y, x)| & two numbers & \verb|rad| \\
\verb|math.exp(x)| & a number & \verb|float| \\
\verb|math.log(x)|, \verb|log2|, \verb|log10| & a number above 0 & \verb|float| \\
\bottomrule
\end{tabular}
\end{center}

\begin{zeno}
import math

print(math.sqrt(2.0), math.pi)
print(math.floor(-1.5), math.ceil(1.2), math.round(2.5), math.round(-2.5))
print(math.log(math.e), math.log2(8.0), math.log10(1000.0), math.exp(0.0))
\end{zeno}

\subsection{Rounding}

\verb|floor| rounds down, \verb|ceil| rounds up, and \verb|round| goes
to the nearest whole number, with a half going away from zero. All
three give an \verb|int|, because an index or a register value is
usually what is wanted next. To show a number with two decimals,
don't round it; format it: \verb|f"{x:.2f}"|.

\begin{why}
Python's \verb|round| sends a half to the nearest even number, so
\verb|round(2.5)| is 2 and \verb|round(3.5)| is 4. That is better for
adding up many rounded numbers, and a surprise to nearly everyone who
meets it. Zeno rounds as school taught, as C and Rust do.
(Decision 6.25.)
\end{why}

\subsection{Logarithms}

\verb|log| is the natural logarithm, with base $e$, as in C, Python,
JavaScript and Java. \verb|log2| and \verb|log10| have their bases in
their names. There is no \verb|log(x, base)|: \verb|log2(8.0)| is
exactly 3, while dividing two natural logarithms gives
\verb|3.0000000000000004|. (Decision 6.26.)

\subsection{A number with no answer}

The square root of a negative number, or the logarithm of 0, has no
answer, and is a \verb|ValueError| with the number in its message:

\begin{zeno}[fails=yes]
import math
try {
    print(math.sqrt(-1.0))
} except ValueError as e {
    print(e.message)
}
print(math.log(0.0))
\end{zeno}

\begin{why}
C gives a \verb|nan|, "not a number", which spreads through every sum
it touches and is printed twenty lines later, far from the line that
made it. An error at the line that asked is easier to find.
(Decision 6.27.)
\end{why}

\section{Angles}

Degrees handed to a function that counts in radians is the classic
mistake with angles, and in C it gives a wrong number with no
complaint. In Zeno an angle can carry its unit
(\hyperref[ch-units]{Numbers with Units}), and the functions of
\verb|math| read it:

\begin{zeno}
import math

print(math.sin(30 deg), math.sin(0.5 rad), math.sin(0.5))
print(math.cos(180 deg), math.tan(45 deg))
\end{zeno}

A plain number is taken too, and means radians, because that is what
the sine of a number means in mathematics, and signal code is full of
\verb|math.sin(2.0 * math.pi * f * t)|. The sine of 30 degrees is not
exactly 0.5, because a float cannot hold $\pi/6$ exactly; Python
prints the same digits.

The functions that give an angle, \verb|asin|, \verb|acos|,
\verb|atan| and \verb|atan2|, give it in \verb|rad|, so the value says
what it is. \verb|.deg| changes it:

\begin{zeno}
import math

let heading = math.atan2(1.0, -1.0)
print(heading, heading.deg)
print(heading / 1 rad)
\end{zeno}

\verb|atan2(y, x)| takes \verb|y| first. It is the angle of the point
$(x, y)$ seen from the middle, and unlike \verb|atan(y / x)| it knows
which side of the circle the point is on.

\section{What is not here}

\begin{center}
\begin{tabular}{ll}
\toprule
Elsewhere & In Zeno \\
\midrule
\verb|pow(x, y)| & \verb|x ** y| \\
\verb|trunc(x)| & \verb|int(x)| \\
\verb|fabs(x)| & \verb|abs(x)| \\
\verb|degrees(x)|, \verb|radians(x)| & \verb|angle.deg|, \verb|angle.rad| \\
\verb|round(x, 2)| & \verb|f"{x:.2f}"| \\
\bottomrule
\end{tabular}
\end{center}

The answers are the C library's, in every way of running a program:
a compiled program calls \verb|sin| and \verb|log| in the C library,
and \verb|zeno run| reaches the same functions. Two machines with
different C libraries can differ in the last digit of an answer that
is not exact. (Decision 6.28.)

\section{Five small instruments}

This program, from the \verb|examples| folder of Zeno, is five pieces
of arithmetic that firmware does every day: a compass from a
magnetometer, a thermistor, a level meter in decibels, a capacitor
charging, and a sine table for an 8-bit converter.

\begin{zeno}[src=examples/instruments.zeno]
\end{zeno}

\begin{summary}
\verb|abs|, \verb|min|, \verb|max| and \verb|sum| need no import.
\verb|math| holds \verb|pi|, \verb|e|, \verb|sqrt|, \verb|floor|,
\verb|ceil|, \verb|round| (half away from zero, giving an int), the
angle functions, \verb|exp| and three logarithms. Angles go in with
their units, or as radians, and come out in \verb|rad|. A number with
no answer is a \verb|ValueError|, never \verb|nan|.
\end{summary}

\begin{tryit}
\begin{enumerate}
\item Write a function \verb|distance| that takes the four
  \verb|float|s \verb|x1|, \verb|y1|, \verb|x2| and \verb|y2|, and
  gives the distance between the two points.
\item Make a sine table of 32 entries for a 10-bit converter (0 to
  1023).
\item Work out how many bits are needed to count to a number with
  \verb|log2| and \verb|ceil|. Check it for 255, 256 and 1000.
\item What does \verb|math.asin(1.5)| do? Catch it and print the
  message.
\end{enumerate}
\end{tryit}
